% Part: many-valued-logic
% Chapter: three-valued-logics
% Section: goedel

\documentclass[../../../include/open-logic-section]{subfiles}

\begin{document}

\olfileid{mvl}{thr}{god}

\olsection{ગોડેલ તર્કશાસ્ત્રો}

કુર્ત ગોડેલે $n$-મૂલ્ય તર્કશાસ્ત્રોની એવી શ્રેણી રજૂ કરી કે
તેમાંના દરેકમાં સ્ફુરણાવાદી તર્કશાસ્ત્રમાં પ્રામાણ્ય ધરાવતાં બધાં
!!{formula}s સામેલ છે, અને તે દરેક તર્કશાસ્ત્ર શાસ્ત્રીય
તર્કશાસ્ત્રમાં સમાયેલું છે.
તેમાંનું પહેલું રસપ્રદ તર્કશાસ્ત્ર આ છે:

\begin{defn}\ollabel{defn:goedel}
\emph{$3$-મૂલ્ય ગોડેલ તર્કશાસ્ત્ર}~$\LogGod$ નીચેની શ્રેણિક વડે
વ્યાખ્યાયિત થાય છે:
\begin{enumerate}
  \item $\lfalse$, $\lnot$, $\land$, $\lor$, $\lif$
  ધરાવતી પ્રમાણભૂત વિધાનાત્મક ભાષા~$\Lang L_0$.
  \item સત્યમૂલ્યોનો ગણ $V = \{\True, \Undef, \False\}$.
  \item એકમાત્ર નિર્દિષ્ટ મૂલ્ય $\True$ છે; એટલે $V^+ = \{\True\}$.
  \item $\lfalse$ માટે $\tf{\lfalse} = \False$ છે. અન્ય
  સંયોજકોનાં સત્યવિધેયો નીચેની સારણીઓ વડે અપાય છે:
  \begin{center}
    \begin{tabular}{c|c}
      $\tf{\lnot}[\LogGod]$ & \\
      \hline
      $\True$ & $\False$ \\
      $\Undef$ & $\False$ \\
      $\False$ & $\True$
    \end{tabular}
    \quad
    \begin{tabular}{c|ccc}
      $\tf{\land}[\LogGod]$ & $\True$ & $\Undef$ & $\False$ \\
      \hline
      $\True$ & $\True$ & $\Undef$ & $\False$ \\
      $\Undef$ & $\Undef$ & $\Undef$ & $\False$\\
      $\False$ & $\False$ & $\False$ & $\False$
    \end{tabular}
    \\[2ex]
    \begin{tabular}{c|ccc}
      $\tf{\lor}[\LogGod]$ & $\True$ & $\Undef$ & $\False$ \\
      \hline
      $\True$ & $\True$ & $\True$ & $\True$ \\
      $\Undef$ & $\True$ & $\Undef$ & $\Undef$ \\
      $\False$ & $\True$ & $\Undef$ & $\False$
    \end{tabular}
    \quad
    \begin{tabular}{c|ccc}
      $\tf{\lif}[\LogGod]$ & $\True$ & $\Undef$ & $\False$ \\
      \hline
      $\True$ & $\True$ & $\Undef$ & $\False$ \\
      $\Undef$ & $\True$ & $\True$ & $\False$ \\
      $\False$ & $\True$ & $\True$ & $\True$
    \end{tabular}
  \end{center}
\end{enumerate}
\end{defn}

અહીં $\land$ અને~$\lor$ની સત્યતાસારણીઓ લુકાશેવિચ અને પ્રબળ
ક્લીની તર્કશાસ્ત્ર જેવી જ છે, પણ $\lnot$ અને~$\lif$ની
સત્યતાસારણીઓ દરેકમાં જુદી છે. ગોડેલ તર્કશાસ્ત્રમાં
$\tf{\lnot}(\Undef) = \False$. લુકાશેવિચ અને ક્લીની
તર્કશાસ્ત્રોથી વિપરીત, $\tf{\lif}(\Undef, \False) = \False$;
અને ક્લીની તર્કશાસ્ત્રથી વિપરીત, પરંતુ લુકાશેવિચ તર્કશાસ્ત્રની જેમ,
$\tf{\lif}(\Undef, \Undef) = \True$.

ઉપર સૂચવેલા સ્ફુરણાવાદી તર્કશાસ્ત્ર સાથેના સંબંધ પરથી સમજાય છે કે
$\LogGod[3]$ તેની નજીક છે. બધી સ્ફુરણાવાદી સત્યતાઓ
$\LogGod[3]$માં પુનરુક્તિ છે; જ્યારે સ્ફુરણાવાદી અર્થમાં પ્રામાણ્ય
ન ધરાવતી ઘણી શાસ્ત્રીય પુનરુક્તિઓ $\LogGod[3]$માં પણ પુનરુક્તિ નથી.
દાખલા તરીકે, આ પુનરુક્તિઓ નથી:
\begin{align*}
  & p \lor \lnot p && (p \lif q) \lif (\lnot p \lor q) \\
  & \lnot\lnot p \lif p && \lnot(\lnot p \land \lnot q) \lif (p \lor q) \\
  & ((p \lif q) \lif p) \lif p && \lnot(p \lif q) \lif (p \land \lnot q)
\end{align*}
પરંતુ $\LogGod[3]$ની દરેક પુનરુક્તિ સ્ફુરણાવાદી અર્થમાં પણ
પ્રામાણ્ય ધરાવે એવું નથી; દાખલા તરીકે $\lnot\lnot p \lor \lnot p$
અથવા $(p \lif q) \lor (q \lif p)$.

\begin{prob}
  નીચેનાં સૂત્રો~$\LogGod[3]$માં પુનરુક્તિ છે તે બતાવવા
  સત્યતાસારણીઓ આપો:
  \begin{align*}
    & \lnot\lnot p \lor \lnot p\\
    & (p \lif q) \lor (q \lif p) \\
    & \lnot(p \land q) \lif (\lnot p \lor \lnot q) \\
    & (p \lif q) \lor (q \lif r) \lor (r \lif s)
  \end{align*}
\end{prob}

\begin{prob}
  નીચેનાં સૂત્રો~$\LogGod[3]$માં પુનરુક્તિ નથી તે બતાવવા
  સત્યતાસારણીઓ આપો:
  \begin{align*}
    & (p \lif q) \lif (\lnot p \lor q) \\
    & \lnot(\lnot p \land \lnot q) \lif (p \lor q) \\
    & ((p \lif q) \lif p) \lif p \\
    & \lnot(p \lif q) \lif (p \land \lnot q)
  \end{align*}
\end{prob}

\begin{prob}
  ગોડેલ તર્કશાસ્ત્રમાં નીચેનામાંથી કયા ફલિતતા-સંબંધો સાચા છે?
  દરેક માટે સત્યતાસારણી આપો.
  \begin{enumerate}
    \item $p, p \lif q \Entails q$
    \item $p \lor q, \lnot p \Entails q$
    \item $p \land q \Entails p$
    \item $p \Entails p \land p$
    \item $p \Entails p \lor q$
  \end{enumerate}
\end{prob}

\end{document}
